\chapter{Fourier Transform of Distributions}
\label{ch:iii19}

The Fourier transform extends to tempered distributions by duality. The same algebraic rules for differentiation, multiplication, translation, and convolution then survive far beyond integrable functions.

\section*{Learning goals}
After this chapter, the reader should be able to:
\begin{itemize}
\item define and compute Fourier transforms of tempered distributions by duality;
\item derive transform formulas for Dirac masses, derivatives, polynomials, and exponentials in distributional form;
\item use differentiation and multiplication identities without assuming classical integrability;
\item solve algebraic multiplier equations in frequency space at the distribution level;
\item check constants and signs by testing transform identities against Schwartz functions;
\item distinguish ordinary-function Fourier formulas from their distributional extensions;
\end{itemize}
\section*{Conceptual roadmap}
\[
\boxed{\text{structure}\;\longrightarrow\;\text{estimate}\;\longrightarrow\;\text{limit or representation}\;\longrightarrow\;\text{application}.}
\]

\paragraph{Mathematical setting.}
Use $\langle T,\phi\rangle=T(\phi)$ with linear tests. Exactly $\langle\widehat T,\phi\rangle=\langle T,\widehat\phi\rangle$ and $\langle\mathcal F^{-1}T,\phi\rangle=\langle T,\mathcal F^{-1}\phi\rangle$. Fubini verifies compatibility with the integral transform for regular integrable functions. Thus $\mathcal F^2T=T(-\,\cdot)$, $\widehat{x_jT}=-(2\pi i)^{-1}D_j\widehat T$, and $\widehat{\delta_a}=e^{-2\pi ia\cdot\xi}$. Continuity of the automorphism holds for both the weak dual and strong dual topologies.

\section{Dual definition}
For $T\in\mathcal S'$, define $\widehat T(\phi)=T(\widehat\phi)$, with the linear pairing and no inverse placement.

\section{Fourier automorphism}
Because the Fourier transform is a topological automorphism of Schwartz space, it induces an automorphism of tempered distributions.

\section{Dirac transforms}
Point masses transform into complex exponentials, while constants transform into Dirac masses under the standard convention.


% BEGIN VOL03-EXPANSION III19-example-01
\begin{example}[The transform of Dirac is the constant distribution]\label{ex:iii19-expand-01}
Under the chapter's \(2\pi\) convention,
\[
\widehat{\delta_0}=1.
\]
Indeed for a Schwartz test function, the dual definition gives the action of
\(\delta_0\) on the transformed test, which is the transform value at zero.
The transform definition at zero is the integral of the original test
function, exactly the regular distribution defined by the constant \(1\).
\end{example}
% END VOL03-EXPANSION III19-example-01
\section{Differentiation}
Distributional differentiation becomes multiplication by polynomial frequency factors.


% BEGIN VOL03-EXPANSION III19-example-02
\begin{example}[Distributional differentiation keeps the classical multiplier rule]\label{ex:iii19-expand-02}
For \(T\in\mathcal S'\),
\[
\widehat{\partial^\alpha T}
=(2\pi i\xi)^\alpha\widehat T.
\]
The proof is dual: move the derivative to the test function, use the classical
Schwartz-space transform identity, then move multiplication back to the
transformed distribution. No pointwise differentiability or integrability of
\(T\) is required.
\end{example}
% END VOL03-EXPANSION III19-example-02
\section{Polynomial multiplication}
Multiplication by $x^\alpha$ becomes differentiation in frequency.

\section{Translation and modulation}
The classical transform identities extend by duality.

\section{Convolution with Schwartz functions}
A tempered distribution can be convolved with a Schwartz function to produce a smooth function with controlled growth.

\section{PDE multipliers}
Constant-coefficient differential equations become algebraic multiplier equations in frequency space.


% BEGIN VOL03-EXPANSION III19-example-03
\begin{example}[Solving a multiplier equation with a singular source]\label{ex:iii19-expand-03}
For the one-dimensional equation
\[
-u''=\delta_0,
\]
Fourier transformation gives
\[
4\pi^2\xi^2\widehat u=1.
\]
Formally,
\[
\widehat u=(4\pi^2\xi^2)^{-1},
\]
which is not locally integrable at zero and is not an ordinary function multiplier. A precise extension is
\[
\widehat u=-\frac1{4\pi^2}D_\xi\operatorname{pv}\frac1\xi.
\]
Here $\langle\operatorname{pv}(1/\xi),\phi\rangle=\lim_{\varepsilon\downarrow0}\int_{|\xi|>\varepsilon}\phi(\xi)/\xi\,d\xi$; the limit exists by subtracting $\phi(0)$ near zero. Since $\xi\operatorname{pv}(1/\xi)=1$, differentiation gives $\xi^2D_\xi\operatorname{pv}(1/\xi)=-1$, verifying the multiplier equation. In physical space $u=-|x|/2$ satisfies $-u''=\delta_0$. Any two tempered solutions differ by an affine function (equivalently by delta and delta-prime terms in frequency). The example shows why
distributional Fourier multipliers are needed precisely where the symbol
vanishes.
\end{example}
% END VOL03-EXPANSION III19-example-03
\section{Core structural results}

\begin{theorem}[Fourier transform of Dirac]\label{thm:iii19-01}
Under the $2\pi$ convention, $\widehat{\delta_0}=1$ as a tempered distribution.
\begin{proof}
For $\phi\in\mathcal S$, $\widehat{\delta_0}(\phi)=\widehat\phi(0)=\int\phi$, directly by the transform definition.
\end{proof}
\end{theorem}

\begin{theorem}[Derivative rule]\label{thm:iii19-02}
$\widehat{\partial^\alpha T}=(2\pi i\xi)^\alpha\widehat T$.
\begin{proof}
Move the derivative onto the test function, use the classical Schwartz-space transform identity, and then return the multiplication to the transformed distribution.
\end{proof}
\end{theorem}

\begin{theorem}[Fourier automorphism]\label{thm:iii19-03}
The Fourier transform is a continuous bijection of $\mathcal S'$ with continuous inverse.
\begin{proof}
Dualize the continuous bijection of Schwartz space and use the inverse Fourier transform for the inverse operator.
\end{proof}
\end{theorem}

\section{Worked examples}

\begin{example}[Transform of the Dirac mass]\label{ex:iii19-worked-01}
Define the Fourier transform on tempered distributions by duality. Then
\[
\langle\widehat{\delta_0},\varphi\rangle
=\langle\delta_0,\widehat\varphi\rangle
=\widehat\varphi(0)
=\int\varphi(x)\,dx.
\]
Therefore \(\widehat{\delta_0}=1\) under the volume convention.
\end{example}

\begin{example}[Transform of the constant distribution]\label{ex:iii19-worked-02}
By the inverse duality calculation,
\[
\widehat{1}=\delta_0
\]
under the symmetric \(2\pi\)-frequency convention used in the volume. The pair
\(1\leftrightarrow\delta_0\) expresses perfect localization duality.
\end{example}

\begin{example}[Distributional derivatives obey the same multiplier rule]\label{ex:iii19-worked-03}
For \(T\in\mathcal S'\),
\[
\widehat{D_jT}=2\pi i\xi_j\,\widehat T.
\]
The proof transfers \(D_j\) to the Schwartz test function, uses the classical
transform identity there, and then returns to the dual pairing.
\end{example}

\begin{example}[Transform of a translated point mass]\label{ex:iii19-worked-04}
For \(\delta_a\),
\[
\widehat{\delta_a}(\xi)=e^{-2\pi ia\cdot\xi}
\]
as a regular tempered distribution. Thus a point translation becomes a pure
frequency phase.
\end{example}

\section{Solved dossiers}
Each dossier combines several ideas from the chapter in a longer argument. The aim is to make hypotheses, calculations, and proof strategy explicit.

\begin{problem}[Transform of the Dirac mass]\label{prob:iii19-dossier-01}
Compute \(\widehat{\delta_0}\) as a tempered distribution.
\end{problem}
\begin{solution}
Duality gives
\(\langle\widehat{\delta_0},\varphi\rangle=\int\varphi\), so the transform is
the constant distribution \(1\).
\end{solution}

\begin{problem}[Transform of the constant distribution]\label{prob:iii19-dossier-02}
Compute the Fourier transform of the constant distribution \(1\).
\end{problem}
\begin{solution}
It is \(\delta_0\) under the declared convention.
\end{solution}

\begin{problem}[Distributional derivatives obey the same multiplier rule]\label{prob:iii19-dossier-03}
Prove the derivative-multiplier identity for tempered distributions.
\end{problem}
\begin{solution}
Apply the definition of distributional derivative and Fourier duality,
then use the classical Schwartz identity
\(\widehat{D_j\varphi}=2\pi i\xi_j\widehat\varphi\).
\end{solution}

\begin{problem}[Transform of a translated point mass]\label{prob:iii19-dossier-04}
Compute \(\widehat{\delta_a}\).
\end{problem}
\begin{solution}
Use translation of \(\delta_0\), or evaluate the dual definition; the
result is \(e^{-2\pi ia\cdot\xi}\).
\end{solution}

\begin{problem}[Laplacian multiplier]\label{prob:iii19-dossier-05}
Compute the Fourier multiplier for \(-\Delta T\).
\end{problem}
\begin{solution}
\(\widehat{-\Delta T}=4\pi^2|\xi|^2\widehat T\).
\end{solution}

\begin{problem}[Multiplication by coordinates]\label{prob:iii19-dossier-06}
How does multiplication by \(x_j\) act after Fourier transform?
\end{problem}
\begin{solution}
Under the convention,
\(\widehat{x_jT}=-(2\pi i)^{-1}\partial_{\xi_j}\widehat T\), interpreted
distributionally.
\end{solution}

\begin{problem}[Transform is an automorphism]\label{prob:iii19-dossier-07}
Why does the Fourier transform map \(\mathcal S'\) to itself continuously?
\end{problem}
\begin{solution}
It is defined by duality from the continuous automorphism
\(\mathcal F:\mathcal S\to\mathcal S\).
\end{solution}

\begin{problem}[PDE motivation]\label{prob:iii19-dossier-08}
Why is the distributional transform useful for constant-coefficient PDE?
\end{problem}
\begin{solution}
Differential operators become polynomial multipliers even when solutions or
forcing terms are singular distributions.
\end{solution}

\begin{problem}[Fourier transform of Dirac proof dossier]\label{prob:iii19-dossier-09}
Prove the following structural statement and identify the step where its main hypothesis is used: Under the $2\pi$ convention, $\widehat{\delta_0}=1$ as a tempered distribution.
\end{problem}
\begin{solution}
For $\phi\in\mathcal S$, $\widehat{\delta_0}(\phi)=\widehat\phi(0)=\int\phi$, directly by the transform definition. This proof also explains why the stated hypothesis is structural rather than cosmetic.
\end{solution}

\begin{problem}[Derivative rule proof dossier]\label{prob:iii19-dossier-10}
Prove the following structural statement and identify the step where its main hypothesis is used: $\widehat{\partial^\alpha T}=(2\pi i\xi)^\alpha\widehat T$.
\end{problem}
\begin{solution}
Move the derivative onto the test function, use the classical Schwartz-space transform identity, and then return the multiplication to the transformed distribution. This proof also explains why the stated hypothesis is structural rather than cosmetic.
\end{solution}

\begin{problem}[Fourier automorphism proof dossier]\label{prob:iii19-dossier-11}
Prove the following structural statement and identify the step where its main hypothesis is used: The Fourier transform is a continuous bijection of $\mathcal S'$ with continuous inverse.
\end{problem}
\begin{solution}
Dualize the continuous bijection of Schwartz space and use the inverse Fourier transform for the inverse operator. This proof also explains why the stated hypothesis is structural rather than cosmetic.
\end{solution}

\begin{problem}[Chapter synthesis]\label{prob:iii19-dossier-12}
Connect the main definitions and structural theorems of this chapter into one proof strategy. Explain what object is controlled, what estimate or closure principle is used, and what conclusion becomes available.
\end{problem}
\begin{solution}
The Fourier transform extends to tempered distributions by duality. The same algebraic rules for differentiation, multiplication, translation, and convolution then survive far beyond integrable functions. The recurring strategy is to encode the object in the chapter's natural measurable, normed, Fourier, or distributional structure, establish the controlling estimate, and then pass to the desired limit or representation.
\end{solution}
\section{Exercises with complete solutions}
The shorter exercise layer remains separate from the solved dossiers.

\begin{exercise}\label{exr:iii19-01}
State and apply the central principle behind Dual definition. Give one immediate consequence in the setting of this chapter.
\end{exercise}
\begin{hint}
Define \(\widehat T\) by duality: pair \(T\) with the Fourier transform of the Schwartz test function.
\end{hint}
\begin{solution}
For $T\in\mathcal S'$, define $\widehat T(\phi)=T(\widehat\phi)$, with the linear pairing and no inverse placement. As an immediate consequence, the same principle can be inserted into later convergence, approximation, transform, or weak-form arguments whenever its hypotheses are verified.
\end{solution}

\begin{exercise}\label{exr:iii19-02}
State and apply the central principle behind Fourier automorphism. Give one immediate consequence in the setting of this chapter.
\end{exercise}
\begin{hint}
To transform a Dirac mass, evaluate the transformed test function at the support point and simplify the exponential.
\end{hint}
\begin{solution}
Because the Fourier transform is a topological automorphism of Schwartz space, it induces an automorphism of tempered distributions. As an immediate consequence, the same principle can be inserted into later convergence, approximation, transform, or weak-form arguments whenever its hypotheses are verified.
\end{solution}

\begin{exercise}\label{exr:iii19-03}
State and apply the central principle behind Dirac transforms. Give one immediate consequence in the setting of this chapter.
\end{exercise}
\begin{hint}
Distributional differentiation becomes multiplication by the frequency variable, with constants determined by the transform convention.
\end{hint}
\begin{solution}
Point masses transform into complex exponentials, while constants transform into Dirac masses under the standard convention. As an immediate consequence, the same principle can be inserted into later convergence, approximation, transform, or weak-form arguments whenever its hypotheses are verified.
\end{solution}

\begin{exercise}\label{exr:iii19-04}
State and apply the central principle behind Differentiation. Give one immediate consequence in the setting of this chapter.
\end{exercise}
\begin{hint}
Multiplication by \(x\) becomes differentiation in frequency; derive it by duality rather than formal pointwise manipulation.
\end{hint}
\begin{solution}
Distributional differentiation becomes multiplication by polynomial frequency factors. As an immediate consequence, the same principle can be inserted into later convergence, approximation, transform, or weak-form arguments whenever its hypotheses are verified.
\end{solution}

\begin{exercise}\label{exr:iii19-05}
State and apply the central principle behind Polynomial multiplication. Give one immediate consequence in the setting of this chapter.
\end{exercise}
\begin{hint}
Polynomials transform into derivatives of Dirac masses, so expect singular frequency-space objects rather than ordinary functions.
\end{hint}
\begin{solution}
Multiplication by $x^\alpha$ becomes differentiation in frequency. As an immediate consequence, the same principle can be inserted into later convergence, approximation, transform, or weak-form arguments whenever its hypotheses are verified.
\end{solution}

\begin{exercise}\label{exr:iii19-06}
State and apply the central principle behind Translation and modulation. Give one immediate consequence in the setting of this chapter.
\end{exercise}
\begin{hint}
To solve a constant-coefficient equation, transform it and divide by the symbol only where that operation is distributionally meaningful.
\end{hint}
\begin{solution}
The classical transform identities extend by duality. As an immediate consequence, the same principle can be inserted into later convergence, approximation, transform, or weak-form arguments whenever its hypotheses are verified.
\end{solution}

\begin{exercise}\label{exr:iii19-07}
State and apply the central principle behind Convolution with Schwartz functions. Give one immediate consequence in the setting of this chapter.
\end{exercise}
\begin{hint}
Check signs and \(2\pi\) factors by applying the identity to one Schwartz test function and tracing the convention.
\end{hint}
\begin{solution}
A tempered distribution can be convolved with a Schwartz function to produce a smooth function with controlled growth. As an immediate consequence, the same principle can be inserted into later convergence, approximation, transform, or weak-form arguments whenever its hypotheses are verified.
\end{solution}

\begin{exercise}\label{exr:iii19-08}
State and apply the central principle behind PDE multipliers. Give one immediate consequence in the setting of this chapter.
\end{exercise}
\begin{hint}
Do not require classical integrability of the original object; tempered-distribution duality is exactly what replaces it.
\end{hint}
\begin{solution}
Constant-coefficient differential equations become algebraic multiplier equations in frequency space. As an immediate consequence, the same principle can be inserted into later convergence, approximation, transform, or weak-form arguments whenever its hypotheses are verified.
\end{solution}


% BEGIN VOL03-EXPANSION III19-exercises-01
\section{Graded supplementary exercises}
These exercises extend the preserved foundational set and are graded by mathematical role.

\subsection*{Standard computations}

\begin{exercise}[Dirac transform]\label{exr:iii19-expand-01}
What is the Fourier transform of delta\_0 under the chapter convention?
\end{exercise}
\begin{hint}
Use the dual definition.
\end{hint}
\begin{solution}
The constant distribution 1.
\end{solution}

\begin{exercise}[Constant transform]\label{exr:iii19-expand-02}
What is the transform of the constant distribution 1?
\end{exercise}
\begin{hint}
Use Fourier duality/inversion.
\end{hint}
\begin{solution}
delta\_0, with the chapter normalization.
\end{solution}

\begin{exercise}[Derivative]\label{exr:iii19-expand-03}
What is the transform of DT?
\end{exercise}
\begin{hint}
Use the derivative rule.
\end{hint}
\begin{solution}
2 pi i xi times hat T.
\end{solution}

\begin{exercise}[Translated Dirac]\label{exr:iii19-expand-04}
What is the transform of delta\_a?
\end{exercise}
\begin{hint}
Use translation of delta.
\end{hint}
\begin{solution}
The exponential exp(-2 pi i a xi).
\end{solution}

\begin{exercise}[Modulated constant]\label{exr:iii19-expand-05}
What distribution transforms to delta at a nonzero frequency?
\end{exercise}
\begin{hint}
Use modulation.
\end{hint}
\begin{solution}
A complex exponential in x.
\end{solution}

\subsection*{Proofs}

\begin{exercise}[Dirac proof]\label{exr:iii19-expand-06}
Prove the transform of delta\_0 is 1.
\end{exercise}
\begin{hint}
Evaluate the dual definition on a Schwartz test.
\end{hint}
\begin{solution}
The transformed test at zero equals the integral of the original test.
\end{solution}

\begin{exercise}[Derivative rule]\label{exr:iii19-expand-07}
Prove the distributional derivative multiplier rule.
\end{exercise}
\begin{hint}
Move derivatives through duality.
\end{hint}
\begin{solution}
The Schwartz transform identity supplies the factor.
\end{solution}

\begin{exercise}[Automorphism]\label{exr:iii19-expand-08}
Why is Fourier transform a bijection on tempered distributions?
\end{exercise}
\begin{hint}
Dualize the Schwartz automorphism.
\end{hint}
\begin{solution}
The inverse Schwartz transform induces the inverse on the dual.
\end{solution}

\begin{exercise}[Translation rule]\label{exr:iii19-expand-09}
Derive the transform of delta\_a.
\end{exercise}
\begin{hint}
Evaluate delta\_a on the transformed test.
\end{hint}
\begin{solution}
It yields the phase exponential multiplying the integral functional.
\end{solution}

\subsection*{Counterexamples and hypothesis testing}

\begin{exercise}[Ordinary formula failure]\label{exr:iii19-expand-10}
Why can hat(delta) not be computed by the ordinary L1 integral formula?
\end{exercise}
\begin{hint}
Delta is not an L1 function.
\end{hint}
\begin{solution}
The dual distribution definition is essential.
\end{solution}

\begin{exercise}[Division by zeros]\label{exr:iii19-expand-11}
Why is dividing by a PDE symbol delicate where the symbol vanishes?
\end{exercise}
\begin{hint}
The reciprocal may not be an ordinary function.
\end{hint}
\begin{solution}
One needs a distributional extension and may gain nonuniqueness from homogeneous solutions.
\end{solution}

\begin{exercise}[Polynomial transform]\label{exr:iii19-expand-12}
Why do transforms of polynomials become derivatives of delta rather than ordinary decaying functions?
\end{exercise}
\begin{hint}
Multiplication by x becomes frequency differentiation.
\end{hint}
\begin{solution}
Starting from transform of 1=delta generates delta derivatives.
\end{solution}

\subsection*{Applications and investigations}

\begin{exercise}[Impulse response]\label{exr:iii19-expand-13}
Why is Fourier transform of delta central to linear systems?
\end{exercise}
\begin{hint}
Delta is the identity for convolution.
\end{hint}
\begin{solution}
Transforming it to 1 makes convolution responses and transfer functions algebraic.
\end{solution}

\begin{exercise}[PDE Green kernel]\label{exr:iii19-expand-14}
How does a distributional multiplier produce a fundamental solution?
\end{exercise}
\begin{hint}
Invert the operator symbol in the distribution sense and inverse transform.
\end{hint}
\begin{solution}
The resulting distribution E satisfies P(D)E=delta.
\end{solution}

\subsection*{Challenge problems}

\begin{exercise}[Transform of delta prime]\label{exr:iii19-expand-15}
Compute the Fourier transform of delta'\_0.
\end{exercise}
\begin{hint}
Apply the derivative rule to delta.
\end{hint}
\begin{solution}
It is 2 pi i xi.
\end{solution}

\begin{exercise}[Transform of x]\label{exr:iii19-expand-16}
Describe the Fourier transform of the regular tempered distribution x.
\end{exercise}
\begin{hint}
Use multiplication by x as frequency differentiation of transform of 1.
\end{hint}
\begin{solution}
$\widehat{x}=-(2\pi i)^{-1}\delta_0'=i\delta_0'/(2\pi)$, by the multiplication rule applied to $\widehat1=\delta_0$.
\end{solution}
% END VOL03-EXPANSION III19-exercises-01
\section*{Chapter summary}
The chapter now supplies a canonical theorem-and-dossier layer for subsequent measure, Fourier, distributional, Sobolev, and PDE arguments.
